% Section 9, part 1 of 4 (continued in parts 2-4); label prefix: gam. No new macros.
% Local letters are declared where they are used. Discs are written \bbD: the letter D is the operator q d/dq.

\section{The Gamma identity and the normalisation}
\label{sec:gam}

This section proves Theorem Gamma (\cref{thm:gam-main}): the full-Stirling saddle series $c(z;L)$ of \cref{def:gam-c} has no pole at $L=1$, that is at $q=0$, and $c(z;1)=h(z)$, the series $h=h_{\bw_0}$ of (QRR). By Theorem~P and the glue of \cref{sec:con} this is equivalent to $\Dn=1$, and it gives $\Fg{g}$ for every $g\geq 2$ (\cref{thm:gam-D}). The proof is local: near a point at which exactly one Gamma factor of the Mellin--Barnes summand degenerates, only that factor is written as an Euler integral; after a change of chart the two-variable stationary phase depends holomorphically on a parameter that vanishes at the point, where a Gaussian computation evaluates it exactly (Theorem~LMI, \cref{thm:gam-lmi}). At $L=0$, that is at $q=\infty$, the same theorem gives the free coefficients of \cref{lem:gau-HPF} in closed form (\cref{thm:gam-ginf}, \cref{cor:gam-free}).

\subsection{The saddle series}
\label{sec:gam-series}

We use $c(z;L)=\SPh_L[\mathsf S_0;A_{\mathrm{St}}]$ and its data as in \cref{def:gam-c}; $c_n\in\AL$ does not depend on the choices of disc and branches (\cref{prop:gam-rational}). At each of the three points $L=1$, $L=\zeta^i$ and $L=0$ exactly one of the six Gamma factors of the summand degenerates, namely $\Gamma((1-\sigma)/z)$, $\Gamma((\zeta^i-\sigma)/z)$ and the twist factor $\Gamma(1+5\sigma/z)$ respectively, and \cref{thm:gam-lmi} applies at each of them; it is used at $L=1$ and $L=0$, and the points $\zeta^i$ are reached by rotation (\cref{thm:gam-main}~(iii)).

\subsection{Local Mellin inversion}
\label{sec:gam-lmi}

Data that are holomorphic at the base point enter only through their Taylor series.

\begin{lem}[The delta lemma with a source]\genusstatus{AI-proposed, unchecked}
\label{lem:gam-deltasrc}
Let $a_*\in\bbC^\times$ and $c\in\bbC$, let $P\in s^2\bbC\fps{s}$, and let $A\in\bbC\fps{s}\fps{z}$. In two variables $(s,m)$ consider the phase $F(s,m):=P(s)+a_*sm$, of exactly this form (no term in $m$ alone and no term $s^km$ with $k\geq 2$), and the amplitude $\re^{cm}A(s,z)$. Then $(0,0)$ is a non-degenerate critical point of $F$, $\det\Hess F(0,0)=-a_*^2$, and, exactly,
\begin{equation}
\label{eq:gam-deltasrc}
\SPn_{(0,0)}\bigl[F;\,\re^{cm}A\bigr]=\exp\bigl(P(\delta)/z\bigr)\,A(\delta,z),\qquad \delta:=-cz/a_*,
\end{equation}
where both factors are re-expanded in $z$: $P(\delta)/z=\sum_{k\geq 2}\coef{s^k}P\,(-c/a_*)^kz^{k-1}$ lies in $z\bbC\fps{z}$, and $A(\delta,z)=\sum_{n\geq 0}z^n\sum_{j+l=n}\coef{s^jz^l}A\,(-c/a_*)^j$. For $c=0$ the series is $A(0,z)$.
\end{lem}

\begin{proof}
The gradient $(P'(s)+a_*m,\,a_*s)$ vanishes at $(0,0)$, and the Hessian there is $H=\bigl(\begin{smallmatrix}P''(0)&a_*\\ a_*&0\end{smallmatrix}\bigr)$, $\det H=-a_*^2\neq0$; the covariance $C=-H^{-1}$ has $C_{ss}=0$, $C_{sm}=-1/a_*$, $C_{mm}=P''(0)/a_*^2$, and the part of order at least three of $F$ is $F_3(s)=P(s)-P''(0)s^2/2$. Write $\re^{F_3(s)/z}A(s,z)=\sum_{a\geq 0}g_a(z)s^a$, with $g_a$ of order at least $-a/3$ in $z$. By \cref{def:fsp-wick} the left side of \eqref{eq:gam-deltasrc} is $\sum_{a,r\geq 0}\frac{c^r}{r!}g_a(z)\langle s^am^r\rangle_Cz^{(a+r)/2}$. As $C_{ss}=0$, in the matching sum of \cref{lem:fsp-ring}~(iii) every copy of $s$ is paired with a copy of $m$; so $\langle s^am^r\rangle_C=\frac{r!}{(r-a)!}(r-a-1)!!\,C_{sm}^{\,a}C_{mm}^{\,(r-a)/2}$ if $r-a\geq0$ is even ($(-1)!!:=1$), and $0$ otherwise. With $r=a+2j$ and $(2j-1)!!/(2j)!=1/(2^jj!)$,
\[
\sum_{r\geq 0}\frac{c^r}{r!}\langle s^am^r\rangle_C\,z^{(a+r)/2}=(cC_{sm}z)^a\exp\Bigl(\frac{c^2C_{mm}z}{2}\Bigr)=\delta^a\exp\Bigl(\frac{P''(0)\delta^2}{2z}\Bigr).
\]
Hence the left side is $\exp(P''(0)\delta^2/(2z))\sum_ag_a(z)\delta^a=\exp(P''(0)\delta^2/(2z))\exp(F_3(\delta)/z)A(\delta,z)$, which is the right side. All sums are finite at each order of $z$, since $g_a(z)\delta^a$ has order at least $2a/3$; and the substitution $s\mapsto\delta$ is compatible with sums, products and exponentials on series $\sum_ag_a(z)s^a$ with $g_a$ of order at least $-a/3$, a class containing $F_3/z$ and $A$ and closed under products.
\end{proof}

\begin{defn}[LMI data]
\label{def:gam-lmidata}
LMI data consist of a point $s_*\in\bbC$, a number $a_*\in\bbC^\times$, an integer $M_*$, an open disc $\bbD$ around $s_*$, a holomorphic function $\varphi$ on $\bbD$, and a series $\cF(\sigma,z)=\sum_{m\geq 0}\cF_m(\sigma)z^m$ of holomorphic functions on $\bbD$, such that $\cF_0$ has no zero on $\bbD$ (in particular $\cF(s_*,0)\neq 0$) and $a_*+(\sigma-s_*)\varphi''(\sigma)\neq 0$ for every $\sigma\in\bbD$. Put $\varphi_k:=\varphi^{(k)}(s_*)/k!$, put
\[
a(\sigma):=a_*(\sigma-s_*),
\]
which is exactly linear in $\sigma$, and $U_*:=\bbD\setminus\{s_*\}$. For $L\in U_*$ choose a branch of $\log a$ near $L$ and put
\[
\lambda(L):=-\varphi'(L)-a_*\log a(L),\qquad S(\sigma):=\lambda(L)\,\sigma+\varphi(\sigma)+a(\sigma)\log a(\sigma)-a(\sigma),
\]
\[
A(\sigma,z):=a(\sigma)^{M_*-1/2}\,\mathrm{Poch}_{M_*}\bigl(z/a(\sigma)\bigr)\exp\Bigl(\sum_{k\geq 1}\bk_k\bigl(z/a(\sigma)\bigr)^{2k-1}\Bigr)\,\cF(\sigma,z),
\]
that is, the expansion of $z^{M_*+a/z}\,\Gst(M_*+a/z)$ given by \cref{lem:fsp-stirling}~(iii) and \cref{lem:fsp-stirfe}, without its factor $(2\pi z)^{1/2}$ and without its phase factor $\re^{(a\log a-a)/z}$, times $\cF$; the power $a^{M_*-1/2}$ is taken with any branch near $L$. Let
\[
c^{\cF}(z;L)=\sum_{n\geq 0}c^{\cF}_n(L)z^n:=\SPh_L\bigl[S;\,A\bigr].
\]
Finally put $\delta:=-M_*z/a_*$, $\tilde\varphi(x):=\varphi(s_*+x)-\varphi(s_*)-\varphi'(s_*)x$, and define the chart parameter
\[
\mathsf Q(L):=\re^{-\lambda(L)/a_*}=a(L)\,\re^{\varphi'(L)/a_*}.
\]
\end{defn}

The point $L$ is critical for $S$: since the derivative of $a\log a-a$ is $a_*\log a$, $S'(\sigma)=\lambda(L)+\varphi'(\sigma)+a_*\log a(\sigma)$, which vanishes at $\sigma=L$. Moreover $S''(L)=\varphi''(L)+a_*^2/a(L)$, that is $S''(L)=(a_*+(L-s_*)\varphi''(L))/(L-s_*)$, which is non-zero on $U_*$; and $A(L,0)=a(L)^{M_*-1/2}\cF_0(L)\neq 0$; so $c^{\cF}(z;L)$ is defined for every $L\in U_*$.

% Section 9, part 2 of 4: Theorem LMI and its proof. Continued in part 3.

\begin{thm}[Theorem LMI: local Mellin inversion]\genusstatus{AI-proposed, unchecked}
\label{thm:gam-lmi}
Let LMI data be given (\cref{def:gam-lmidata}).
\begin{enumerate}
\item For $L\in U_*$ every coefficient $c^{\cF}_n(L)$ depends only on $L$, not on the branch of $\log a$ or on $\lambda(L)$, and $c^{\cF}_n$ is holomorphic on $U_*$.
\item Every $c^{\cF}_n$ extends holomorphically to $L=s_*$.
\item With these extensions,
\begin{equation}
\label{eq:gam-lmi}
c^{\cF}(z;s_*)=\exp\bigl(\tilde\varphi(\delta)/z\bigr)\,\frac{\cF(s_*+\delta,z)}{\cF(s_*,0)},\qquad \delta=-\frac{M_*z}{a_*},
\end{equation}
where both factors are re-expanded in $z$: $\tilde\varphi(\delta)/z=\sum_{k\geq 2}\varphi_k(-M_*/a_*)^kz^{k-1}$ and $\cF(s_*+\delta,z)=\sum_{n\geq 0}z^n\sum_{m+j=n}\cF_m^{(j)}(s_*)\,(-M_*/a_*)^j/j!$.
\end{enumerate}
\end{thm}

That $a$ is exactly linear in $\sigma$ and that $\varphi$ and $\cF$ do not involve $\lambda$ makes the chart data of Step~2 of the proof holomorphic at $L=s_*$; the condition $\cF(s_*,0)\neq 0$ is used for the normalisation at $s_*$.

\begin{proof}
\emph{Step 0: (i).} For $k\geq 2$ the Taylor coefficient of order $k$ at $L$ of $a\log a-a$ is $(-1)^ka_*/(k(k-1)(L-s_*)^{k-1})$; it involves no branch, and $\lambda(L)$ enters only the term of order one of $S$, which vanishes at the critical point. The normalised amplitude $A(L+t,z)/A(L,0)$ is the product of the binomial series $(1+t/(L-s_*))^{M_*-1/2}$, of the series $\mathrm{Poch}_{M_*}(z/a(L+t))\exp(\sum_k\bk_k(z/a(L+t))^{2k-1})$ in $t$, where $a(L+t)=a(L)(1+t/(L-s_*))$, and of $\cF(L+t,z)/\cF_0(L)$; it involves neither $\lambda$ nor a branch. By \cref{lem:fsp-univ}, $c^{\cF}_n(L)$ is a polynomial with rational coefficients in finitely many of these Taylor coefficients and in $-1/S''(L)=-(L-s_*)/(a_*+(L-s_*)\varphi''(L))$. Each of them is a holomorphic function of $L$ on $U_*$; the Taylor coefficients of $\varphi$, and those of the $\cF_m$ divided by $\cF_0(L)$, are holomorphic on all of $\bbD$, where $\cF_0$ has no zero. This proves (i).

\emph{Step 1: two variables.} Fix $L\in U_*$ and the branch of $\log a$ near $L$ used in $\lambda=\lambda(L)$. We use a local letter $\chi$ for the logarithm of the variable of the Euler integral of $\Gamma(M_*+a/z)$, and consider, in the two variables $(\sigma,\chi)$, the formal integrand with phase and amplitude
\[
\Phi_2(\sigma,\chi):=\lambda\sigma+\varphi(\sigma)+a(\sigma)\chi-\re^{\chi},\qquad \re^{M_*\chi}\,\cF(\sigma,z),
\]
at the point $(L,\chi_L)$, $\chi_L:=\log a(L)$ (the chosen branch). It is a formal integrand in the sense of \cref{sec:fsp}, motivated by Euler's integral $\int_0^\infty x^{M_*+a/z-1}\re^{-x/z}\,\dd x=z^{M_*+a/z}\,\Gamma(M_*+a/z)$, $a=a(\sigma)$, in the variable $x=\re^\chi$. The gradient $(\lambda+\varphi'(\sigma)+a_*\chi,\ a(\sigma)-\re^{\chi})$ vanishes at $(L,\chi_L)$ by the choice of $\lambda$, and, since $a''=0$, the Hessian there is
\[
\begin{pmatrix}\varphi''(L)&a_*\\ a_*&-a(L)\end{pmatrix},\qquad\text{with determinant }-a(L)\varphi''(L)-a_*^2=-a(L)S''(L)\neq 0.
\]
We apply \cref{lem:fsp-fubini} with the inner variable $\chi$ and the outer variable $\sigma$. The inner Hessian $-a(L)$ is invertible. Over the coefficient ring $\bbC\fps{\sigma-L}$ the inner critical point is $\chi(\sigma):=\chi_L+\log(1+(\sigma-L)/(L-s_*))$, the branch of $\log a(\sigma)$ that continues $\chi_L$, with $\re^{\chi(\sigma)}=a(\sigma)$; the inner critical value $\lambda\sigma+\varphi(\sigma)+a(\sigma)\chi(\sigma)-a(\sigma)$ is $S(\sigma)$; the inner Hessian at $\chi(\sigma)$ is $-a(\sigma)$, so that the factor of \cref{lem:fsp-fubini} is $(a(\sigma)/a(L))^{-1/2}$. The inner series is computed by \cref{lem:fsp-cov} with the map $\chi\mapsto x:=\re^{\chi}$, which sends $\chi(\sigma)$ to $a(\sigma)$ and has the Jacobian $a(\sigma)$ there; it turns the phase $a(\sigma)\log x-x$ into $a(\sigma)\chi-\re^{\chi}$ and the amplitude $x^{M_*-1}$ into $a(\sigma)^{-1}\re^{M_*\chi}$. With \cref{lem:fsp-stirling} for the unit $a=a(\sigma)$ of $\bbC\fps{\sigma-L}$, the integer $M=M_*$ and the measure $\dd x$, this gives
\begin{align*}
\SPn_{\chi(\sigma)}\bigl[a(\sigma)\chi-\re^{\chi};\,\re^{M_*\chi}\bigr]&=a(\sigma)\,\SPn_{a(\sigma)}\bigl[a(\sigma)\log x-x;\,x^{M_*-1}\bigr]\\
&=a(\sigma)^{M_*}\,\mathrm{Poch}_{M_*}\bigl(z/a(\sigma)\bigr)\exp\Bigl(\sum_{k\geq 1}\bk_k\bigl(z/a(\sigma)\bigr)^{2k-1}\Bigr).
\end{align*}
The factor $\cF(\sigma,z)\in\bbC\fps{\sigma-L}\fps{z}$ does not involve $\chi$ and is carried along by the linearity of the inner $\SPn$ over $\bbC\fps{\sigma-L}\fps{z}$ (\cref{def:fsp-wick}). So the outer amplitude is
\[
\Bigl(\frac{a(\sigma)}{a(L)}\Bigr)^{-1/2}a(\sigma)^{M_*}\,\mathrm{Poch}_{M_*}\bigl(z/a(\sigma)\bigr)\exp\Bigl(\sum_{k\geq1}\bk_k\bigl(z/a(\sigma)\bigr)^{2k-1}\Bigr)\cF(\sigma,z)=a(L)^{1/2}A(\sigma,z),
\]
and \cref{lem:fsp-fubini} gives
\begin{equation}
\label{eq:gam-step1}
\SPh_{(L,\chi_L)}\bigl[\Phi_2;\,\re^{M_*\chi}\cF\bigr]=\SPh_L\bigl[S;\,A\bigr]=c^{\cF}(z;L).
\end{equation}

\emph{Step 2: the chart.} Let $\mu$ be a second local letter and translate $\chi=\mu-\lambda/a_*$ (\cref{lem:fsp-cov}, Jacobian $1$). Put $\mu_c(L):=-\varphi'(L)/a_*$ and $p(L):=(L,\mu_c(L))$. The translation maps $p(L)$ to $(L,\chi_L)$, since $\mu_c(L)-\lambda/a_*=-\varphi'(L)/a_*+\varphi'(L)/a_*+\log a(L)=\chi_L$. Because $a$ is exactly $a_*(\sigma-s_*)$, we have $\lambda\sigma+a(\sigma)(\mu-\lambda/a_*)=\lambda s_*+a(\sigma)\mu$, and the translation turns the phase into
\[
\Phi_2(\sigma,\mu-\lambda/a_*)=\lambda s_*+\Psi_{\mathsf Q}(\sigma,\mu),\qquad \Psi_{\mathsf Q}(\sigma,\mu):=\varphi(\sigma)+a(\sigma)\mu-\mathsf Q\,\re^{\mu},
\]
where $\Psi_{\mathsf Q}$ is a phase local to this proof (as is $\Psi_0$ below), and the amplitude into $\re^{M_*(\mu-\lambda/a_*)}\cF=\mathsf Q^{M_*}\re^{M_*\mu}\cF$, with $\mathsf Q=\mathsf Q(L)=\re^{-\lambda/a_*}$ and $\mathsf Q^{M_*}:=\re^{-M_*\lambda/a_*}$. The constant $\lambda s_*$ of the phase changes only the critical value (\cref{lem:fsp-const}), and the constant factor $\mathsf Q^{M_*}$ of the amplitude cancels in $\SPh$. With \eqref{eq:gam-step1} we obtain, for $L\in U_*$,
\begin{equation}
\label{eq:gam-chart}
c^{\cF}(z;L)=\SPh_{p(L)}\bigl[\Psi_{\mathsf Q(L)};\,\re^{M_*\mu}\cF\bigr].
\end{equation}
Here $\lambda$ has disappeared, except through $\mathsf Q(L)=a(L)\re^{\varphi'(L)/a_*}$. This function is single-valued: a change of the branch of $\log a(L)$ changes $\lambda$ by an integer multiple of $2\pi\ri a_*$ and leaves $\re^{-\lambda/a_*}$ unchanged. It is holomorphic on $\bbD$, and it has a simple zero at $s_*$.

\emph{The family of critical points.} For every $L\in\bbD$, including $L=s_*$, the point $p(L)$ is a critical point of $\Psi_{\mathsf Q(L)}$: the derivative $\partial_\sigma\Psi=\varphi'(\sigma)+a_*\mu$ vanishes at $p(L)$, and $\partial_\mu\Psi=a(\sigma)-\mathsf Q(L)\re^{\mu}$ equals $a(L)-a(L)\re^{\varphi'(L)/a_*}\re^{-\varphi'(L)/a_*}=0$ there. Since $\mathsf Q(L)\re^{\mu_c(L)}=a(L)$, the Hessian at $p(L)$ is
\[
H(L)=\begin{pmatrix}\varphi''(L)&a_*\\ a_*&-a(L)\end{pmatrix},\qquad \det H(L)=-a_*\bigl(a_*+(L-s_*)\varphi''(L)\bigr)\neq 0\quad(L\in\bbD).
\]
At $L=s_*$ we have $\mathsf Q(s_*)=0$, $\Psi_0(\sigma,\mu)=\varphi(\sigma)+a(\sigma)\mu$ and $p(s_*)=(s_*,\mu_*)$ with $\mu_*:=-\varphi'(s_*)/a_*$.

\emph{Holomorphy: (ii).} In the variables $s=\sigma-L$ and $m=\mu-\mu_c(L)$ the phase $\Psi_{\mathsf Q(L)}$ is
\[
\Psi_{\mathsf Q(L)}(p(L))+\bigl[\varphi(L+s)-\varphi(L)-\varphi'(L)s\bigr]+a_*sm-a(L)\bigl(\re^{m}-1-m\bigr),
\]
because the terms of order one cancel; so its Taylor coefficients are the $\varphi^{(k)}(L)/k!$, the number $a_*$ and the numbers $-a(L)/k!$. The normalised amplitude is $\re^{M_*m}\cF(L+s,z)/\cF_0(L)$, with the Taylor coefficients $M_*^j/j!$ and $\cF_l^{(j)}(L)/(j!\,\cF_0(L))$; and the entries of $-H(L)^{-1}$ are holomorphic on $\bbD$, because $\det H(L)$ has no zero there. By \cref{lem:fsp-univ} the $n$-th coefficient of the right side of \eqref{eq:gam-chart} is a polynomial with rational coefficients in finitely many of these functions (compare \cref{lem:fsp-param}~(i)). It is therefore holomorphic on all of $\bbD$, and by \eqref{eq:gam-chart} it equals $c^{\cF}_n$ on $U_*$. This proves (ii). At $L=s_*$ the data are the Taylor data of $\Psi_0$ and of $\re^{M_*\mu}\cF$ at $(s_*,\mu_*)$ and the entries of $-H(s_*)^{-1}$; since the coefficients are given by the same universal polynomials,
\begin{equation}
\label{eq:gam-atQ0}
c^{\cF}(z;s_*)=\SPh_{(s_*,\mu_*)}\bigl[\Psi_0;\,\re^{M_*\mu}\cF\bigr].
\end{equation}

\emph{Step 3: the value.} In the variables $s:=\sigma-s_*$ and $m:=\mu-\mu_*$,
\[
\Psi_0=\varphi(s_*)+\tilde\varphi(s)+a_*sm
\]
exactly: the terms $\varphi'(s_*)s+a_*\mu_*s$ cancel, and there is no term in $m$ alone and no term $s^km$ with $k\geq 2$. The amplitude is $\re^{M_*\mu_*}\re^{M_*m}\cF(s_*+s,z)$. The constant $\varphi(s_*)$ of the phase does not matter (\cref{lem:fsp-const}), and \cref{lem:gam-deltasrc} with $P=\tilde\varphi\in s^2\bbC\fps{s}$, $c=M_*$ and $A(s,z)=\cF(s_*+s,z)$ gives $\SPn_{(s_*,\mu_*)}[\Psi_0;\,\re^{M_*\mu}\cF]=\re^{M_*\mu_*}\exp(\tilde\varphi(\delta)/z)\cF(s_*+\delta,z)$, whose constant term is $\re^{M_*\mu_*}\cF(s_*,0)$. Dividing by it, \eqref{eq:gam-atQ0} gives \eqref{eq:gam-lmi}.
\end{proof}

% Section 9, part 3 of 4: Theorem Gamma, the hatted series, the value of D. Continued in part 4.

\subsection{Theorem Gamma and the value of \texorpdfstring{$\Dn$}{D}}
\label{sec:gam-thm}

Recall $h=h_{\bw_0}$ from \eqref{eq:set-hw}, with $\bw_0=\{1-\zeta^j\}_{j=1,\dots,4}\cup\{-5\}$, and $h_n=\coef{z^n}h$. Put
\[
\mathsf T_k:=5^{1-2k}+\sum_{j=1}^4(\zeta^j-1)^{1-2k}\qquad(k\geq 1).
\]
Since $1-2k$ is odd, $\mathsf T_k=-\sum_{a\in\bw_0}a^{1-2k}$, and therefore $h(z)=\exp\bigl(\sum_{k\geq 1}\bk_k\mathsf T_kz^{2k-1}\bigr)$.

\begin{thm}[Theorem Gamma]\genusstatus{AI-proposed, unchecked}
\label{thm:gam-main}
Let $c(z;L)$ be the saddle series of \cref{def:gam-c}.
\begin{enumerate}
\item Every $c_n$ lies in $\AL$.
\item No $c_n$ has a pole at $L=1$, that is at $q=0$, and
\begin{equation}
\tag{$\Gamma$}
\label{eq:gam-Gamma}
c(z;1)=h(z),\qquad\text{that is, }c_n(1)=h_n\ \text{for every }n\geq 0.
\end{equation}
\item $c_n(\zeta L)=\zeta^{-n}c_n(L)$ for every $n$. Hence no $c_n$ has a pole at a fifth root of unity, $c_n$ lies in $\bbQ[L,L^{-1}]$, every exponent of $L$ that occurs in $c_n$ is congruent to $4n$ modulo $5$, and $c(z;\zeta^i)=h_{\bw_i}(z)=h(\zeta^{-i}z)$.
\item No $c_n$ has a pole at $L=0$, that is at $q=\infty$, and $c(z;0)=\Ginf(z)$, where
\[
\Ginf(z):=\exp\Bigl(\sum_{l\geq 1}(-1)^{l+1}\frac{B_{5l+1}(1/5)}{l(5l+1)}\,z^{5l}\Bigr),
\]
with $B_k(\theta)$ the Bernoulli polynomials; this is the series displayed on the right of (37) of \cite[Lemma~7.7]{GJR18} at $m=5$, $k=1$ and equivariant parameter $1$ in the notation there, read as a series in $z$ (the coefficients on the left of (37) are introduced in \cite[Corollary~7.6]{GJR18} through powers of $-z$); it is obtained there for their Landau--Ginzburg integral; the value $c(z;0)=\Ginf(z)$ is proved here independently (\cref{thm:gam-ginf}), and the comparison with \cite{GJR18} enters no proof. Consequently $c_n\in\bbQ[L]$ (\cref{cor:gam-poly}).
\end{enumerate}
\end{thm}

Only (ii) is used in the proof of the main theorem; (iii) and (iv) are by-products, and (iv) identifies the free coefficients (\cref{cor:gam-free}).

\begin{proof}
(i) is \cref{prop:gam-rational}. (iv) is proved in \cref{thm:gam-ginf} and \cref{cor:gam-poly} below; for the statements of (iv) these proofs use only (i), (iii) (and through it (ii)), \cref{thm:gam-lmi} and \cref{lem:fsp-dlmf5118} (the further statements of \cref{cor:gam-poly} also use \cref{thm:con-Pa}~(iv)).

(ii) We apply \cref{thm:gam-lmi} with $s_*=1$, $a_*=-1$, so that $a(\sigma)=1-\sigma$, and $M_*=0$. Let $\bbD$ be a disc around $1$ that contains none of $0,\zeta,\dots,\zeta^4$, fix branches of $\log(5\sigma)$ and of $\log(\zeta^j-\sigma)$, $j=1,\dots,4$, on $\bbD$, and put
\[
\varphi(\sigma):=5\sigma\log(5\sigma)-5\sigma+\sum_{j=1}^4\bigl[(\zeta^j-\sigma)\log(\zeta^j-\sigma)-(\zeta^j-\sigma)\bigr],
\]
\[
\cF(\sigma,z):=(5\sigma)^{1/2}\prod_{j=1}^4(\zeta^j-\sigma)^{-1/2}\exp\Bigl(\sum_{k\geq 1}\bk_kz^{2k-1}\Bigl[(5\sigma)^{1-2k}+\sum_{j=1}^4(\zeta^j-\sigma)^{1-2k}\Bigr]\Bigr),
\]
the phase and the amplitude of the five factors of the summand that do not degenerate at $\sigma=1$. Thus $\cF$ is $A_{\mathrm{St}}$ without its factor $(1-\sigma)^{-1/2}\exp(\sum_k\bk_kz^{2k-1}(1-\sigma)^{1-2k})$ of index $j=0$. The functions $\varphi$ and $\cF_m$ are holomorphic on $\bbD$, $\cF_0(1)=5^{1/2}\prod_{j\geq 1}(\zeta^j-1)^{-1/2}\neq 0$, and after shrinking $\bbD$ the function $\cF_0$ has no zero on $\bbD$ and $-1+(\sigma-1)\varphi''(\sigma)\neq 0$ there. The term $a\log a-a=(1-\sigma)\log(1-\sigma)-(1-\sigma)$ of \cref{def:gam-lmidata} is the term $j=0$ of $\mathsf S_0$, and, as $\mathrm{Poch}_0=1$, the Stirling form of $\Gamma(0+a/z)$ contributes $a^{-1/2}\exp(\sum_k\bk_k(z/a)^{2k-1})$, the factor $j=0$ of $A_{\mathrm{St}}$. Hence $S=\mathsf S_0$ and $A=A_{\mathrm{St}}$, and $\lambda(L)=-\varphi'(L)+\log(1-L)=-5\log(5L)+\sum_{j=0}^4\log(\zeta^j-L)$ is the logarithm of \cref{def:gam-c}, for the branches chosen. Therefore $c^{\cF}(z;L)=c(z;L)$ for $L\in U_*$. By \cref{thm:gam-lmi}~(ii) every $c_n$ extends holomorphically to $L=1$; since $c_n$ is rational by (i), it has no pole at $L=1$, and its value there is that of the extension. By \eqref{eq:gam-lmi} with $\delta=0$, and since the factor $(5\sigma)^{1/2}\prod_{j\geq 1}(\zeta^j-\sigma)^{-1/2}$ of $\cF$ does not involve $z$,
\[
c(z;1)=\frac{\cF(1,z)}{\cF(1,0)}=\exp\Bigl(\sum_{k\geq 1}\bk_k\mathsf T_kz^{2k-1}\Bigr)=h(z).
\]

(iii) Consider the Taylor data of $\mathsf S_0$, $\mathsf S_1$, $\mathsf S_{2k}$ at $L$ displayed after \cref{def:gam-c}. Since $p_k(\zeta L)=\zeta^{-k}p_k(L)$, the Taylor coefficients at $\zeta L$ are $\zeta^{-k}$, $\zeta^{-k}$ and $\zeta^{1-2k-r}$ times those at $L$. Substitute $t=\zeta t'$ (\cref{lem:fsp-cov}, a linear map with constant Jacobian). In the variable $t'$ the phase $\mathsf S_0(\zeta L+\zeta t')-\mathsf S_0(\zeta L)$ is $\zeta$ times $\mathsf S_0(L+t')-\mathsf S_0(L)$, the normalised factor $\exp(\mathsf S_1(\zeta L+\zeta t')-\mathsf S_1(\zeta L))$ equals $\exp(\mathsf S_1(L+t')-\mathsf S_1(L))$, and $z^{2k-1}\mathsf S_{2k}(\zeta L+\zeta t')=(\zeta^{-1}z)^{2k-1}\mathsf S_{2k}(L+t')$. So the normalised integrand at $\zeta L$ is that at $L$, with the phase multiplied by $\zeta$ and with $z$ replaced by $\zeta^{-1}z$ in the amplitude. \Cref{lem:fsp-homog} with the unit $\zeta$ gives $c(z;\zeta L)=c(\zeta^{-1}z;L)$, that is $c_n(\zeta L)=\zeta^{-n}c_n(L)$. Iterating, $c_n(L)=\zeta^{-in}c_n(\zeta^{-i}L)$; by (ii) the right side has no pole at $L=\zeta^i$, and its value there is $\zeta^{-in}h_n$. So $c(z;\zeta^i)=h(\zeta^{-i}z)$, which is $h_{\bw_i}(z)$ by \cref{lem:set-hrot}. By (i) the only possible poles of $c_n$ are $0$ and the fifth roots of unity; since there are none at the latter, $c_n\in\bbQ[L,L^{-1}]$. Writing $c_n=\sum_ea_eL^e$, the relation $c_n(\zeta L)=\zeta^{-n}c_n(L)$ gives $\zeta^e=\zeta^{-n}$ whenever $a_e\neq 0$, that is $e\equiv -n\equiv 4n$ modulo $5$.
\end{proof}

\begin{prop}[Flat and hatted saddle series]\genusstatus{AI-proposed, unchecked}
\label{prop:gam-hatflat}
For every $n$, $c_n(L)=L^{4n}\,\mathrm{ph}^{\mathrm{SP}}_n(L^{-5})$ as rational functions of $L$. Hence $\mathrm{ph}^{\mathrm{SP}}_n(1)=c_n(1)$ for every $n$, and
\[
\PhiSP(\zh,1)=h(\zh).
\]
\end{prop}

\begin{proof}
The identity of rational functions is \cref{thm:con-Pa}~(iv), which states that every $\mathrm{ph}^{\mathrm{SP}}_n$ is a polynomial in $u$. At $L=1$ it gives $\mathrm{ph}^{\mathrm{SP}}_n(1)=c_n(1)$, and $c_n(1)=h_n$ by \cref{thm:gam-main}~(ii); so $\PhiSP(\zh,1)=\sum_n\mathrm{ph}^{\mathrm{SP}}_n(1)\zh^n=h(\zh)$. The proof of \cref{thm:con-Pa} uses neither this proposition nor \cref{thm:gam-main}.
\end{proof}

\begin{cor}\genusstatus{AI-proposed, unchecked}
\label{cor:gam-equiv}
$\PhiSP(z,1)=h(z)$ holds if and only if the identity \eqref{eq:gam-Gamma} holds. With \cref{cor:con-A2}: $\Dn=1$ if and only if \eqref{eq:gam-Gamma} holds.
\end{cor}

\begin{proof}
By \cref{thm:con-Pa}~(iv), which uses neither \cref{thm:gam-main} nor \eqref{eq:gam-Gamma}, every $c_n$ is a Laurent polynomial in $L$ with $c_n(1)=\mathrm{ph}^{\mathrm{SP}}_n(1)$; so $c(z;1)=\PhiSP(z,1)$, and the first equivalence follows. The second is \cref{cor:con-A2}.
\end{proof}

\begin{thm}[Theorem C: the normalisation is canonical]\genusstatus{AI-proposed, unchecked}
\label{thm:gam-D}
In every conifold frame (\cref{def:set-frames}), $\Dn_0=1$ in $\bbQ\fps{z}$, and $\Dn=1$. Consequently $\Fg{g}$ holds for every $g\geq 2$.
\end{thm}

\begin{proof}
Fix a conifold frame. By \cref{thm:con-glue}, together with \cref{thm:con-Pb}, $\Dn_0(z)=\PhiSP(z,1)/h(z)$ exactly in $\bbQ\fps{z}$; the right side does not depend on the frame (\cref{def:con-PhiSP}). By \cref{prop:gam-hatflat}, which rests on \cref{thm:gam-main}~(ii), $\PhiSP(z,1)=h(z)$. Hence $\Dn_0=1$. By \cref{cor:gau-colours}, $\Dn_i(z)=\Dn_0(\zeta^{-i}z)=1$ for every colour $i$, so $\Dn=\diag(\Dn_i)=1$. Finally $\Fg{g}$, that is $\Dn_0=1$ modulo $z^{3g-2}$ (\cref{def:gau-Fg}), holds for every $g\geq 2$.
\end{proof}

The input (QRR) enters the proof of \cref{thm:gam-D} through \cref{thm:con-glue}, at colour $0$, and through \cref{cor:gau-colours}, whose proof uses the colour covariance of the rows (\cref{cor:gau-rowsNN}~(iii), resting on \cref{lem:gau-colrows}), that is (QRR) for the colours $i\neq0$; \cref{thm:con-glue} also uses the rows and their WKB form, which rest on (M), (G0) and the flatness part of (QRR) (\cref{def:set-rows}~(b)). The proof uses the $R$-matrix of the formal quintic, identified by (T) (\cref{thm:set-GT}), only through \cref{def:set-rows}~(b), $R|_{z=0}=1$ and (QRR).

% Section 9, part 4 of 4: the orbifold point, the free coefficients, prior art. Last part of Section 9.

\begin{thm}[The orbifold point]\genusstatus{AI-proposed, unchecked}
\label{thm:gam-ginf}
No $c_n$ has a pole at $L=0$, and
\begin{equation}
\label{eq:gam-ginf}
c(z;0)=\Ginf(z)=\exp\Bigl(\sum_{l\geq 1}(-1)^{l+1}\frac{B_{5l+1}(1/5)}{l(5l+1)}\,z^{5l}\Bigr).
\end{equation}
\end{thm}

We do not use, and do not claim, that the series of \cite{GJR18} coincides with $c(z;L)$.

\begin{proof}
Apply \cref{thm:gam-lmi} with $s_*=0$, $a_*=5$, so that $a(\sigma)=5\sigma$, and $M_*=1$; then $\delta=-z/5$, and $\mathsf Q=\re^{-\lambda/5}$ is a fifth root of $(-q)^{-1}$. Near $\sigma=0$ only the twist factor $\Gamma(1+5\sigma/z)$ degenerates; the five factors $\Gamma((\zeta^j-\sigma)/z)$, $j=0,\dots,4$, are regular there and supply, on a disc $\bbD$ around $0$ that contains no $\zeta^j$ and with fixed branches,
\[
\varphi(\sigma):=\sum_{j=0}^4\varphi_j(\sigma),\qquad \varphi_j(\sigma):=(\zeta^j-\sigma)\log(\zeta^j-\sigma)-(\zeta^j-\sigma),
\]
\[
\cF:=\prod_{j=0}^4\cF^{(j)},\qquad \cF^{(j)}(\sigma,z):=(\zeta^j-\sigma)^{-1/2}\exp\Bigl(\sum_{k\geq 1}\bk_kz^{2k-1}(\zeta^j-\sigma)^{1-2k}\Bigr).
\]
After shrinking $\bbD$ the conditions of \cref{def:gam-lmidata} hold; $\cF(0,0)=\prod_j(\zeta^j)^{-1/2}\neq 0$. The phase $a\log a-a=5\sigma\log(5\sigma)-5\sigma$ and the Stirling form $a^{1/2}\mathrm{Poch}_1(z/a)\exp(\sum_k\bk_k(z/a)^{2k-1})$ of $\Gamma(1+a/z)$, with $\mathrm{Poch}_1=1$, are the twist terms of $\mathsf S_0$ and of $A_{\mathrm{St}}$; so $S=\mathsf S_0$, $A=A_{\mathrm{St}}$, and $\lambda(L)=-\varphi'(L)-5\log(5L)=\sum_j\log(\zeta^j-L)-5\log(5L)$ is the logarithm of \cref{def:gam-c}. Hence $c^{\cF}(z;L)=c(z;L)$ on $U_*$; by \cref{thm:gam-lmi}~(ii) every $c_n$ extends holomorphically to $L=0$, and, being rational by \cref{thm:gam-main}~(i), has no pole there. By \eqref{eq:gam-lmi},
\[
c(z;0)=\exp\bigl(\tilde\varphi(\delta)/z\bigr)\frac{\cF(\delta,z)}{\cF(0,0)}=\prod_{j=0}^4V_j(z),\qquad V_j(z):=\exp\bigl(\tilde\varphi_j(\delta)/z\bigr)\frac{\cF^{(j)}(\delta,z)}{\cF^{(j)}(0,0)},
\]
with $\tilde\varphi_j(x):=\varphi_j(x)-\varphi_j(0)-\varphi_j'(0)x$. Fix $j$, let $w$ be a local letter with $1/w=\zeta^{-j}z$, and put $\theta:=1/5$. Then $\zeta^j-\delta=z(w+\theta)$, and, expanding $\log(\zeta^j-\delta)=\log\zeta^j+\log(1+\theta/w)$ at $\sigma=0$ (the terms with $\log\zeta^j$ cancel, since $\varphi_j'(0)=-\log\zeta^j$),
\[
\tilde\varphi_j(\delta)/z=(w+\theta)\log(1+\theta/w)-\theta,\qquad \frac{\cF^{(j)}(\delta,z)}{\cF^{(j)}(0,0)}=(1+\theta/w)^{-1/2}\exp\Bigl(\sum_{k\geq 1}\bk_k(w+\theta)^{1-2k}\Bigr),
\]
the second because $z^{2k-1}(z(w+\theta))^{1-2k}=(w+\theta)^{1-2k}$. Therefore
\[
\log V_j=\Bigl(w+\theta-\frac12\Bigr)\log\Bigl(1+\frac{\theta}{w}\Bigr)-\theta+\sum_{k\geq 1}\bk_k(w+\theta)^{1-2k}=\sum_{n\geq 1}\beta_n(\theta)\,w^{-n}
\]
by \cref{lem:fsp-dlmf5118}, the formal shifted Stirling series (no asymptotic statement is used); this is an identity in $\bbQ[\theta]\fps{1/w}$, with the coefficients $\beta_n(\theta)=(-1)^{n+1}B_{n+1}(\theta)/(n(n+1))$ of \cref{def:fsp-stirling}, and the substitution of $1/5$ for $\theta$ and of $\zeta^{-j}z$ for $1/w$ is a ring homomorphism from $\bbQ[\theta]\fps{1/w}$ to $\bbQ(\zeta)\fps{z}$. Summing over $j$,
\[
\log c(z;0)=\sum_{n\geq 1}\beta_n(1/5)\,z^n\sum_{j=0}^4\zeta^{-jn}=5\sum_{l\geq 1}\beta_{5l}(1/5)\,z^{5l},
\]
because $\sum_j\zeta^{-jn}$ is $5$ if $5\mid n$ and $0$ otherwise; and $5\beta_{5l}(1/5)=(-1)^{l+1}B_{5l+1}(1/5)/(l(5l+1))$, since $(-1)^{5l+1}=(-1)^{l+1}$. This is \eqref{eq:gam-ginf}.
\end{proof}

\begin{cor}[The coefficients are polynomials]\genusstatus{AI-proposed, unchecked}
\label{cor:gam-poly}
Every $c_n$ lies in $\bbQ[L]$, and
\[
c_n(L)=\sum_{0\leq k\leq 4n/5}a_{n,k}L^{4n-5k},\qquad a_{n,k}=\coef{u^k}\mathrm{ph}^{\mathrm{SP}}_n.
\]
In particular $\mathrm{ph}^{\mathrm{SP}}_n$ has degree at most $\lfloor 4n/5\rfloor$ in $u$, and $\coef{u^{4m}}\mathrm{ph}^{\mathrm{SP}}_{5m}=c_{5m}(0)$ for $m\geq 1$.
\end{cor}

\begin{proof}
By \cref{thm:gam-main}~(i) and (iii), $c_n\in\bbQ[L,L^{-1}]$, and by \cref{thm:gam-ginf} it has no pole at $L=0$; so $c_n\in\bbQ[L]$. By \cref{prop:gam-hatflat}, that is \cref{thm:con-Pa}~(iv), $c_n(L)=\sum_{k\geq 0}\coef{u^k}\mathrm{ph}^{\mathrm{SP}}_n\,L^{4n-5k}$, a finite sum since $\mathrm{ph}^{\mathrm{SP}}_n$ is a polynomial. The exponents $4n-5k$ are distinct, so $c_n\in\bbQ[L]$ forces $\coef{u^k}\mathrm{ph}^{\mathrm{SP}}_n=0$ whenever $4n-5k<0$, and the coefficient of $L^0$ in $c_{5m}$ is $\coef{u^{4m}}\mathrm{ph}^{\mathrm{SP}}_{5m}$.
\end{proof}

\begin{cor}[The free coefficients]\genusstatus{AI-proposed, unchecked}
\label{cor:gam-free}
For $m\geq 1$ put $\fc^{\mathrm{SP}}_m:=\coef{u^{4m}}\mathrm{ph}^{\mathrm{SP}}_{5m}$, and let $\fc^{\mathrm{can}}_m$ and $\fc^{\mathrm{act}}_m$ be the free coefficients (\cref{lem:gau-HPF}) of $\Phican$ and of $\Phiact$; those of $\Phiact$ are defined because $\Phiact=\Phican$ has coefficients in $\bbQ\fps{u}$ (by \cref{thm:gam-D}, see the proof; here $\Phiact$ is expanded through the representatives $A^{\mathrm{act}}_n\in\bbQ(\zeta)[L,L^{-1}]$ of \cref{thm:gau-R0}~(iii)). Then
\[
\fc^{\mathrm{SP}}_m=c_{5m}(0)=\coef{z^{5m}}\Ginf,\qquad \fc^{\mathrm{can}}_m=\fc^{\mathrm{SP}}_m,\qquad 1+\sum_{m\geq 1}\fc^{\mathrm{can}}_mz^{5m}=\Ginf(z),\qquad \fc^{\mathrm{act}}_m=\fc^{\mathrm{can}}_m.
\]
For instance $\fc^{\mathrm{can}}_1=\coef{z^5}\Ginf=B_6(1/5)/6=4537/3937500$.
\end{cor}

\begin{proof}
\Cref{cor:gam-poly} and \cref{thm:gam-ginf} give $\fc^{\mathrm{SP}}_m=c_{5m}(0)=\coef{z^{5m}}\Ginf$. By \cref{thm:con-Pb}, $\Phican=\PhiSP$, so $\fc^{\mathrm{can}}_m=\fc^{\mathrm{SP}}_m$; this uses neither (QRR) nor \cref{thm:gam-D}. By \cref{thm:gau-R0}~(iv), $\Phiact(\zh,u)\Dn_0(s^4\zh)=\Phican(\zh,u)$, and $\Dn_0=1$ by \cref{thm:gam-D}; so $\Phiact=\Phican$ and $\fc^{\mathrm{act}}_m=\fc^{\mathrm{can}}_m$. The term $l=1$ of \eqref{eq:gam-ginf} is $B_6(1/5)z^5/6$, and $B_6(\theta)=\theta^6-3\theta^5+\tfrac52\theta^4-\tfrac12\theta^2+\tfrac1{42}$ gives $B_6(1/5)/6=4537/3937500$.
\end{proof}

\subsection{Prior art}
\label{sec:gam-rmk}

\begin{rmk}[Prior art]
\label{rmk:gam-prior}
The mechanism (an Euler integral for the degenerate factor, a chart in which the critical point survives at $q=0$, and a factorisation into one-variable Euler integrals there) is Givental's for the equivariant projective spaces \cite[\S10]{MR1901075}, with the chart $x_0=Q/(x_1\cdots x_n)$ and the Stirling series of the factors from \cite[10.4~Lemma]{MR1901075}; the continuity in $Q$ of the asymptotic solution is asserted there. \draftnote{Describe \cite[\S10]{MR1901075} neutrally: decide whether it is presented there as an outline or as a proof.} \cite[Remark~7.2]{GJR18} states the analogous identity for their Landau--Ginzburg integral without proof, and \cite[Lemma~7.7, (37)]{GJR18} gives its value at the orbifold point (\cref{thm:gam-ginf}). See also \cite[Lemma~4.20]{MR2553377} (an exact expansion in $q$ of an oscillatory integral in such a chart) and \cite[\S7, (11)]{MR2276766} (the factor of an Euler-class twist as the stationary-phase expansion of an Euler integral). We are not aware of a proof of \eqref{eq:gam-Gamma} in the literature.
\end{rmk}



